\section{M\'etodo e indicadores comuns}

O artigo adota estrat\'egia descritiva-comparativa, com triangula\c{c}\~ao de fontes oficiais nacionais e internacionais. Para comparabilidade, cada indicador deve registrar fonte, per\'iodo, unidade e abrang\^encia.

\begin{table*}[t]
\centering
\caption{Indicadores comuns para petr\'oleo, g\'as natural e carv\~ao mineral}
\begin{tabular}{p{3.2cm}p{4.0cm}p{2.5cm}p{5.4cm}}
\toprule
Indicador & Unidade de refer\^encia & Periodicidade & Crit\'erio de comparabilidade \\
\midrule
Reservas provadas & barris, bcm ou Mt & anual & convers\~ao para equival\^encia energ\'etica quando necess\'ario \\
Produ\c{c}\~ao & bbl/dia, bcm/ano ou Mt/ano & anual/mensal & explicitar m\'etrica f\'isica e, se usado, equivalente energ\'etico \\
Consumo (demanda) & bbl/dia, bcm/ano ou Mt/ano & anual/mensal & manter distin\c{c}\~ao entre consumo final e demanda prim\'aria \\
Com\'ercio internacional & exporta\c{c}\~oes/importa\c{c}\~oes (volume e valor) & anual & separar balan\c{c}a f\'isica da monet\'aria \\
Pre\c{c}os de refer\^encia & US\$/bbl, US\$/MMBtu, US\$/t & mensal/anual & informar pre\c{c}o nominal, moeda e deflator quando aplic\'avel \\
\bottomrule
\end{tabular}
\end{table*}

Per\'iodo-base comum priorizado: 2015--2025 para s\'eries hist\'oricas, com extens\~ao condicionada \`a disponibilidade e consist\^encia das fontes.
